\section{Step 44 Check capacity before the first write}\label{pcm-wav-s44}
Compare the unsigned destination capacity \texttt{s3} with the canonical total \texttt{s4}. Return status two through \texttt{encode\_capacity} if it is smaller. Perform this before any PCM read or destination write. A top-bit-set capacity represents a large unsigned space, not a negative failure. After this check, retain accepted rate in \texttt{s5}, channels in \texttt{s6} and PCM byte length in \texttt{fp}.

Author region \texttt{e\_capacity}. Use \texttt{s44-1} and stop at \texttt{after\_e\_capacity}.

\paragraph{Observe} A capacity of sixty passes for the trace; fifty-nine fails with unchanged destination and zero count.

\paragraph{Check} How should a capacity with its top bit set compare with sixty?
\begin{itemize}
  \item[A.] As a request for zero bytes.
  \item[B.] As a larger unsigned capacity.
  \item[C.] As a negative capacity.
\end{itemize}
